\chapter{Sums of Squares}

\section{Fermat's Two-Square Theorem}

\begin{theorem}[Fermat's Two-Square Theorem]
  A natural number $n$ is presentable as the sum of two squares if and only if its square-free part is a product of primes congruent to $1$ modulo $4$.
\end{theorem}

\section{Lagrange's Four-Square Theorem}

\begin{theorem}[Lagrange's Four-Square Theorem]
  Every natural number can be written as the sum of four squares.
\end{theorem}

\section{Legendre's Three-Square Theorem}

Characterizing precisely those numbers that can be written as a sum of three squares is harder than the cases for two or four squares, since the numbers expressible as sums of three squares are not closed under multiplication. For example, $3=1^2+1^2+1^2$ and $5=2^2+1^2+0^2$, while $15$ cannot be written as a sum of three squares.

\begin{theorem}[Legendre's Three-Square Theorem]
  A natural number $n$ is presentable as the sum of three squares if and only if it is not of the form $4^k m$ where $m\equiv 7\pmod{8}$.
\end{theorem}


%%% Local Variables:
%%% mode: latex
%%% TeX-master: "elementary-number-theory"
%%% End:

